\documentclass[11pt]{article}
\usepackage[notes]{la}
\lamodule{1}{2}{Row Reduction and Echelon Forms}
\laxdoc{1.1}   % cross-references to Section 1.1 as \ref{1.1-<label>}

\begin{document}

\section{Row Reduction and Echelon Forms}

\begin{guide}
Elimination can end in many different-looking matrices. What makes a matrix ``simple enough'' that
the answer can be read off from it? And how can we organize the work so that it succeeds on
\emph{every} system, including those with no solution or with infinitely many?
\end{guide}

\textbf{By the end of this section you should be able to:}
\begin{itemize}
  \item decide whether a matrix is in echelon form or reduced echelon form;
  \item locate the pivot positions and pivot columns of a matrix;
  \item carry out the row reduction algorithm on any matrix;
  \item write the general solution of a consistent system in parametric form;
  \item use the pivot structure to answer the existence and uniqueness questions, even before the
        system is completely solved.
\end{itemize}

% ---------------------------------------------------------------------
\subsection*{Echelon forms}

In a nonzero row, the \emph{leading entry} is the leftmost nonzero entry. A \emph{zero row} has all entries
equal to $0$.

\begin{defn}{Echelon form and reduced echelon form}{echelon}
A rectangular matrix is in \emph{echelon form} (or \emph{row echelon form}) if
\begin{enumerate}[itemsep=1pt]
  \item every nonzero row lies above every zero row;
  \item the leading entry of each row lies in a column strictly to the right of the leading entry of
        the row above it;
  \item all entries in a column below a leading entry are zero (this follows from~2).
\end{enumerate}
It is in \emph{reduced echelon form} if, in addition,
\begin{enumerate}[itemsep=1pt,start=4]
  \item every leading entry is $1$;
  \item each leading $1$ is the only nonzero entry in its column.
\end{enumerate}
\end{defn}

The leading entries form a ``staircase'' that descends to the right. Schematically (here $\lead$ is any
nonzero number and $\st$ is any number, zero included):
\[
\text{echelon: }
\begin{bmatrix}\lead&\st&\st&\st&\st\\0&\lead&\st&\st&\st\\0&0&0&\lead&\st\\0&0&0&0&0\end{bmatrix}
\qquad
\text{reduced echelon: }
\begin{bmatrix}1&0&\st&0&\st\\0&1&\st&0&\st\\0&0&0&1&\st\\0&0&0&0&0\end{bmatrix}
\]

\begin{kc}{Classify the matrix}{classify}
For each matrix decide: not in echelon form, echelon but not reduced, or reduced echelon. Give a reason.
\[
\text{(a)}\begin{bmatrix}1&2&0&3\\0&0&1&4\\0&0&0&0\end{bmatrix}\quad
\text{(b)}\begin{bmatrix}2&1&5\\0&3&1\\0&0&0\end{bmatrix}\quad
\text{(c)}\begin{bmatrix}1&0&0\\0&0&1\\0&1&0\end{bmatrix}
\]
\[
\text{(d)}\begin{bmatrix}1&3&0&2\\0&1&0&5\\0&0&0&1\end{bmatrix}\quad
\text{(e)}\begin{bmatrix}0&0\\1&0\end{bmatrix}
\]
\end{kc}

A matrix can usually be reduced to many different echelon forms. For instance, with
$A=\begin{bmatrix}1&2\\3&4\end{bmatrix}$, the operation $R_2-3R_1\to R_2$ gives
$\begin{bmatrix}1&2\\0&-2\end{bmatrix}$, which is in echelon form; scaling row~2 by $-\tfrac12$ then gives
$\begin{bmatrix}1&2\\0&1\end{bmatrix}$, also in echelon form. But \emph{reduced} echelon form is different:

\begin{thm}{Uniqueness of the reduced echelon form}{rref}
Every matrix is row equivalent to one and only one matrix in reduced echelon form.
\end{thm}

We accept this theorem here; a proof is given in Appendix~A of the textbook. It is what justifies
speaking of \emph{the} reduced echelon form of $A$ (often abbreviated RREF, as in calculators and
software). In the example above it is $\begin{bmatrix}1&0\\0&1\end{bmatrix}$, no matter how you get there.

\begin{defn}{Pivot position and pivot column}{pivot}
A \emph{pivot position} in a matrix $A$ is a location in $A$ that corresponds to a leading $1$ in the reduced
echelon form of $A$. A \emph{pivot column} is a column of $A$ that contains a pivot position. A \emph{pivot}
is a nonzero number in a pivot position that is used to create zeros below (or above) it.
\end{defn}

\begin{think}{What is well defined?}{pivotpos}
Different row operations lead to different echelon forms, yet we speak of \emph{the} pivot positions of $A$.
Use Theorem~\ref{thm:rref} to explain why this is legitimate. (Hint: what happens to the \emph{positions} of
leading entries when you pass from an echelon form to the reduced echelon form?)
Also: in the example above, the pivots used were $1$ and $-2$. Are these numbers the same as the entries of
$A$ in the pivot positions? Why or why not?
\end{think}

% ---------------------------------------------------------------------
\subsection*{The row reduction algorithm}

The procedure below works on \emph{any} matrix, whether or not it came from a linear system. The
first four steps (the \emph{forward phase}) produce an echelon form; the fifth step (the
\emph{backward phase}) produces the reduced echelon form.

\begin{proc}{Row reduction}{rowreduce}
\begin{enumerate}[itemsep=2pt]
  \item Begin with the leftmost nonzero column. This is a pivot column; the pivot position is at its top.
  \item Select a nonzero entry in this column and, if necessary, interchange rows to move it into the pivot
        position.
  \item Use replacement operations to create zeros in all positions \emph{below} the pivot.
  \item Cover (ignore) the row containing the pivot and all rows above it. Repeat steps~1--3 on the
        submatrix that remains, until there are no nonzero rows left to modify.
  \item \emph{(Backward phase.)} Starting with the rightmost pivot and moving up and to the left, create
        zeros \emph{above} each pivot. Scale each pivot row so that its pivot becomes $1$.
\end{enumerate}
\end{proc}

\begin{exmp}{Row reducing a system with a missing pivot}{rr}
Solve
\[
\begin{aligned}
\phantom{x_1+{}}2x_2+4x_3-2x_4&=6\\
x_1+2x_2+\phantom{4}x_3\phantom{{}-2x_4}&=3\\
2x_1+4x_2+2x_3+\phantom{2}x_4&=5\\
x_1+3x_2+3x_3-\phantom{2}x_4&=6.
\end{aligned}
\]

\textbf{Forward phase.} The leftmost column is a pivot column, but the top entry is $0$, so interchange
rows~1 and~2. Then clear the first column below the pivot, and the second column below the
second pivot (which is the $2$ in row~2):
\begin{align*}
\begin{amat}{cccc|c}0&2&4&-2&6\\1&2&1&0&3\\2&4&2&1&5\\1&3&3&-1&6\end{amat}
&\rops{R_1\leftrightarrow R_2}
\begin{amat}{cccc|c}1&2&1&0&3\\0&2&4&-2&6\\2&4&2&1&5\\1&3&3&-1&6\end{amat}
\rops{R_3-2R_1,\ R_4-R_1}
\begin{amat}{cccc|c}1&2&1&0&3\\0&2&4&-2&6\\0&0&0&1&-1\\0&1&2&-1&3\end{amat}\\[4pt]
&\rops{R_4-\frac12R_2}
\begin{amat}{cccc|c}1&2&1&0&3\\0&2&4&-2&6\\0&0&0&1&-1\\0&0&0&0&0\end{amat}.
\end{align*}
Now look at column~3. Below row~2 it contains only zeros, so there is \emph{nothing to pivot on} there; we
move right to column~4, whose pivot is the $1$ in row~3. The matrix is in echelon form, with pivot columns
$1$, $2$, and $4$. Column~3 is \emph{not} a pivot column.

\textbf{Backward phase.} Clear above the pivot in column~4, then scale row~2, then clear above the pivot in column~2:
\begin{align*}
&\rops{R_2+2R_3}
\begin{amat}{cccc|c}1&2&1&0&3\\0&2&4&0&4\\0&0&0&1&-1\\0&0&0&0&0\end{amat}
\rops{\frac12R_2}
\begin{amat}{cccc|c}1&2&1&0&3\\0&1&2&0&2\\0&0&0&1&-1\\0&0&0&0&0\end{amat}\\[4pt]
&\rops{R_1-2R_2}
\begin{amat}{cccc|c}1&0&-3&0&-1\\0&1&2&0&2\\0&0&0&1&-1\\0&0&0&0&0\end{amat}.
\end{align*}
This is the reduced echelon form. The next subsection shows how to read the solution from it.
\end{exmp}

\begin{remark}[Practical tips]
\begin{itemize}
  \item You may choose which nonzero entry to use as the pivot. Picking an entry that makes the arithmetic
        avoid fractions (for instance a $1$ or $-1$) saves effort and reduces mistakes; interchanging rows
        for this reason is perfectly fine.
  \item Computer programs usually choose the entry of \emph{largest absolute value} in the column as the
        pivot (``partial pivoting''), because dividing by tiny numbers magnifies rounding errors in floating
        point arithmetic.
  \item For an $n\times(n+1)$ augmented matrix the forward phase takes roughly $\tfrac23n^3$ arithmetic
        operations, while the backward phase takes at most about $n^2$. So for large systems the forward
        phase is where nearly all the work is.
  \item Historically, elimination appears in ancient Chinese mathematical texts; it is named after Gauss,
        who used it in the 19th century, and Wilhelm Jordan, who popularized the full reduction to
        reduced echelon form.
\end{itemize}
\end{remark}

% ---------------------------------------------------------------------
\subsection*{Reading the solution from the reduced echelon form}

Take the reduced echelon form from Example~\ref{ex:rr}. It is the augmented matrix of the system
\[
x_1-3x_3=-1,\qquad x_2+2x_3=2,\qquad x_4=-1,\qquad 0=0.
\]
The last equation $0=0$ is always true and imposes no condition.

\begin{defn}{Basic and free variables; general solution}{basicfree}
The variables corresponding to \emph{pivot columns} are the \emph{basic variables}; the others are the
\emph{free variables}. When a system is consistent, we solve each equation for its one basic variable
in terms of the free variables. The result, in which the free variables act as \emph{parameters},
is the \emph{general solution} (or a \emph{parametric description} of the solution set).
\end{defn}

Here the pivot columns are $1,2,4$, so $x_1,x_2,x_4$ are basic and $x_3$ is free:
\[
x_1=-1+3x_3,\qquad x_2=2-2x_3,\qquad x_3\ \text{is free},\qquad x_4=-1.
\]
``$x_3$ is free'' means that \emph{any} real number may be chosen for $x_3$, after which the other
variables are determined. Taking $x_3=0$ gives $(-1,2,0,-1)$; taking $x_3=1$ gives $(2,0,1,-1)$. Check
the second one in the original system:
$0+4+2=6$, $\ 2+0+1=3$, $\ 4+0+2-1=5$, $\ 2+0+3+1=6$. \ $\checkmark$

\begin{kc}{Another member of the solution set}{another}
Use the general solution above to find the solution with $x_3=2$, and verify it in the original
four equations.
\end{kc}

\begin{warn}[Be precise about what is free]
A common error is to write a ``solution'' such as
$x_1=5-x_3,\ x_2=2+x_3,\ x_3=3-x_2$. This describes $x_2$ in terms of $x_3$ \emph{and} $x_3$ in terms of
$x_2$, so it hides which variable is actually free (there is only one free variable in this case).
By convention, \emph{the free variables are the parameters}, and every other variable is expressed in
terms of them alone.
\end{warn}

Notice also what happens with an inconsistent system. The solution set is empty, \emph{even if the
system has free variables}. There is no parametric description of the empty set.

% ---------------------------------------------------------------------
\subsection*{Existence and uniqueness}

The reduced echelon form answers both fundamental questions from Section~1.1.

\begin{thm}{Existence and Uniqueness}{existuniq}
A linear system is consistent if and only if the rightmost column of its augmented matrix is \emph{not}
a pivot column; equivalently, if and only if an echelon form of the augmented matrix has no row of the form
\[
[\,0\ \ \cdots\ \ 0\mid b\,]\qquad\text{with } b\neq0 .
\]
If the system is consistent, then it has
\begin{enumerate}[label=(\roman*),itemsep=1pt]
  \item a \emph{unique} solution when there are no free variables, and
  \item \emph{infinitely many} solutions when there is at least one free variable.
\end{enumerate}
\end{thm}

\begin{proof}
By Theorem~\ref{1.1-thm:rowsol} we may study the reduced echelon form instead. If the rightmost column is a pivot
column, some row of the reduced echelon form is $[\,0\ \cdots\ 0\mid1\,]$, i.e.\ the equation $0=1$, so the
system is inconsistent. (A row $[\,0\ \cdots\ 0\mid b\,]$ with $b\ne0$ in \emph{any} echelon form gives the same
conclusion, and creates a pivot in the last column.)

Conversely, suppose the rightmost column is not a pivot column. Then the leading $1$ of every nonzero row lies
in a coefficient column, so each nonzero equation has the form
\[
(\text{one basic variable})+(\text{terms involving free variables only})=\text{constant}.
\]
Give the free variables any values (say all $0$) and solve each equation for its basic variable: this is a
solution, so the system is consistent. If there are no free variables, the values of the basic variables are forced,
so the solution is unique. If there is at least one free variable, different values of that free variable give
different solutions, so there are infinitely many.
\end{proof}

\begin{proc}{Solving a linear system}{solve}
\begin{enumerate}[itemsep=1pt]
  \item Write the augmented matrix of the system.
  \item Row reduce to an echelon form. If the rightmost column is a pivot column, stop: the system is inconsistent.
  \item Otherwise continue to the reduced echelon form.
  \item Write the system of equations corresponding to the reduced matrix.
  \item Solve each nonzero equation for its basic variable in terms of the free variables.
  \item Check your answer in the \emph{original} system.
\end{enumerate}
\end{proc}

\begin{kc}{Read the pivots}{readpivots}
Each matrix below is the augmented matrix of a system in an echelon form. For each, decide whether the system is
consistent and, if so, whether the solution is unique.
\[
\text{(i)}\begin{amat}{ccc|c}1&2&0&4\\0&0&1&3\end{amat}\qquad
\text{(ii)}\begin{amat}{cc|c}1&0&5\\0&1&2\\0&0&0\end{amat}\qquad
\text{(iii)}\begin{amat}{cc|c}2&1&3\\0&0&4\end{amat}
\]
\end{kc}

% ---------------------------------------------------------------------
\subsection*{What the shape of a matrix tells us}

The theorem lets us predict a great deal from the \emph{pattern of pivots} alone. Let $A$ be the $m\times n$
coefficient matrix of a system with $m$ equations and $n$ unknowns, and write $[A\mid b]$ for its augmented
matrix, where $b$ is the column of right-hand sides.

\begin{thm}{Counting pivots}{counting}
\begin{enumerate}[label=(\alph*),itemsep=2pt]
  \item If $A$ has a pivot position in \emph{every row}, then the system is consistent for \emph{every} choice of $b$.
  \item If $A$ has a pivot position in \emph{every column}, then the system has \emph{at most one} solution,
        for every choice of $b$.
  \item If $m<n$ and the system is consistent, then it has infinitely many solutions.
  \item If $m>n$, then there is some choice of $b$ for which the system is inconsistent.
\end{enumerate}
\end{thm}

\begin{proof}
Row operations on $[A\mid b]$ act on the columns of $A$ exactly as they would on $A$ alone, so the pivot positions
within the first $n$ columns are the pivot positions of $A$.

(a) If every row of $A$ has a pivot, then an echelon form of $A$ has no zero row. So no row of the form
$[\,0\ \cdots\ 0\mid b'\,]$ with $b'\ne0$ can appear, and the system is consistent by
Theorem~\ref{thm:existuniq}.

(b) If every column of $A$ is a pivot column, there are no free variables, so by Theorem~\ref{thm:existuniq} there is
either no solution or exactly one.

(c) There are at most $m$ pivots (one per row) and $m<n$, so some column of $A$ has no pivot, giving a free variable.
If the system is consistent, Theorem~\ref{thm:existuniq} gives infinitely many solutions.

(d) There are at most $n<m$ pivots, so some row of an echelon form $U$ of $A$ is zero. Choose $d$ with last entry $1$;
then $[U\mid d]$ is inconsistent. Since row operations are reversible, applying the reverse operations to $[U\mid d]$
produces some $[A\mid b]$, which is therefore inconsistent as well.
\end{proof}

\begin{kc}{Using the counting theorem}{counting}
(a) $A$ is a $3\times5$ matrix with three pivot positions. Is the system $[A\mid b]$ consistent for every $b$?
When it is consistent, how many solutions does it have?\\
(b) $A$ is a $5\times3$ matrix with three pivot positions. Is the system consistent for every $b$? When it is
consistent, how many solutions does it have?
\end{kc}

\begin{think}{Shape is not destiny}{shape}
\begin{enumerate}[label=(\arabic*)]
  \item Parts (c) and (d) do \emph{not} say that a system with fewer equations than unknowns is consistent, nor that
        a system with more equations than unknowns is inconsistent. Give an example of each of the opposite
        behaviors: an inconsistent system with $m<n$, and a consistent system with $m>n$.
  \item Suppose $A$ is square, $n\times n$. Explain why ``a pivot in every row'' and ``a pivot in every
        column'' are the \emph{same} condition. What do parts (a) and (b) then say together?
        (This innocent-looking observation will grow into one of the central ideas of the course.)
\end{enumerate}
\end{think}

\subsection*{Summary}
\begin{itemize}
  \item \textbf{Echelon form} has a descending staircase of leading entries; \textbf{reduced echelon form} also has
        leading $1$s that are the only nonzero entries in their columns.
  \item Echelon forms are not unique, but the \textbf{reduced echelon form is}, so \textbf{pivot positions} are
        well defined.
  \item \textbf{Forward phase} (to echelon form) answers \emph{existence}; with the \textbf{backward phase} (to
        reduced echelon form) you can read off the solution.
  \item Basic variables $\leftrightarrow$ pivot columns; free variables $\leftrightarrow$ non-pivot columns. A free
        variable gives infinitely many solutions \emph{if the system is consistent}.
  \item A pivot in every row means ``always solvable''; a pivot in every column means ``never more than one solution.''
\end{itemize}

\subsection*{Questions to think about before the next section}
\begin{enumerate}
  \item In your own words: why is the reduced echelon form unique, while an echelon form is not, and why does
        the uniqueness not affect the pivot \emph{positions}?
  \item The augmented matrix of a system of $3$ equations in $4$ unknowns has reduced echelon form with pivots only
        in columns $1$ and $3$. Is it consistent? How many free variables are there?
  \item True or false (justify): ``A system with a free variable is always consistent.''
\end{enumerate}

\subsection*{Answers and discussion}
\begin{description}[leftmargin=0pt,itemsep=5pt]
  \item[Knowledge Check \ref{kc:classify}.] (a)~reduced echelon. (b)~echelon but not reduced: the leading
        entries $2$ and $3$ are not $1$. (c)~not echelon: the leading entry of row~3 (column~2) lies to the
        \emph{left} of the leading entry of row~2 (column~3). (d)~echelon but not reduced: the entries above the
        leading $1$s should be $0$ (the $3$ above the leading $1$ in column~2, and the $2$ and $5$ above the
        leading $1$ in column~4). (e)~not echelon: a zero row lies above a nonzero row.
  \item[Knowledge Check \ref{kc:another}.] With $x_3=2$: $x_1=-1+6=5$, $x_2=2-4=-2$, $x_4=-1$, i.e.\
        $(5,-2,2,-1)$. Check: $-4+8+2=6$, \ $5-4+2=3$, \ $10-8+4-1=5$, \ $5-6+6+1=6$. \ $\checkmark$
  \item[Knowledge Check \ref{kc:readpivots}.] (i)~consistent (last column is not a pivot column); $x_2$ is free,
        so infinitely many solutions: $x_1=4-2x_2$, $x_3=3$. (ii)~consistent with no free variables, so the unique
        solution is $(5,2)$ (the zero row is just $0=0$). (iii)~inconsistent: row~2 says $0=4$.
  \item[Knowledge Check \ref{kc:counting}.] (a)~Yes, by part~(a) of the theorem. There are $5-3=2$ free variables,
        so infinitely many solutions. (b)~No: only $3$ pivots among $5$ rows means some row of an echelon form of $A$
        is zero, so some $b$ gives an inconsistent system. When consistent, the solution is unique, because all
        three columns are pivot columns.
  \item[Pause and think \ref{think:pivotpos}.] \emph{Why ``the'' pivot positions make sense.} Let $U$ be
        \emph{any} echelon form of $A$. To pass from $U$ to the reduced echelon form we only scale rows and add
        multiples of a pivot row to rows \emph{above} it. These steps change the numbers in the matrix but never
        move a leading entry: each stays in its own row and column (and becomes $1$). So the leading entries of
        $U$ sit exactly where the leading $1$s of the reduced echelon form sit. By Theorem~\ref{thm:rref} that
        reduced form is the same matrix no matter which echelon form we started from or which operations we used.
        Hence every echelon form of $A$ has its leading entries in the same positions, and these are the
        pivot positions of $A$.

        \emph{Pivots versus entries.} No. A pivot is the number that sits in the pivot position \emph{after}
        earlier row operations have changed it. For $A=\begin{bmatrix}1&2\\3&4\end{bmatrix}$ the second pivot
        position holds $4$ in $A$, but after $R_2-3R_1\to R_2$ it holds $4-3\cdot2=-2$. If we then scaled row~2
        by $-\tfrac12$, the pivot would become $1$. The \emph{values} of pivots depend on the path taken; the
        \emph{positions} do not.
  \item[Pause and think \ref{think:shape}.]
        \begin{enumerate}[label=(\arabic*),nosep,leftmargin=*]
          \item An inconsistent system with $m<n$: $x_1+x_2+x_3=1,\ x_1+x_2+x_3=2$ (two equations, three
                unknowns; the planes are parallel). A consistent system with $m>n$:
                $x_1=1,\ 2x_1=2,\ 3x_1=3$ (three equations, one unknown); another is
                $x_1+x_2=2,\ x_1-x_2=0,\ 2x_1=2$ from Section~1.1, whose solution is $(1,1)$.
                Extra equations may simply repeat information, and too few equations may still
                contradict each other. Parts (c) and (d) say only that \emph{if} $m<n$ and the system is
                consistent then there are infinitely many solutions, and that when $m>n$ \emph{some} choice
                of $b$ fails, not that every choice fails.
          \item Each pivot lies in a different row and in a different column (that is the staircase).
                Suppose an $n\times n$ matrix has a pivot in every row. Then there are $n$ pivots in $n$
                different columns, but there are only $n$ columns, so every column has a pivot.
                Conversely, if every column has a pivot, there are $n$ pivots in $n$ different rows, so every
                row has a pivot. The two conditions are equivalent. Then (a) says the system is consistent
                for \emph{every} $b$ and (b) says it has \emph{at most one} solution, so together: for every
                $b$ the system has \emph{exactly one} solution. This is the first glimpse of what it means for
                a square matrix to be ``invertible.''
        \end{enumerate}
  \item[Before the next section.] (1)~An echelon form still leaves freedom: you may scale a row by any nonzero
        number, or add multiples of a lower pivot row to a higher row, and each choice gives a different echelon
        form. The extra conditions for reduced echelon form (leading entries equal to $1$, zeros above and below
        each leading $1$) remove exactly that freedom, which is why only one matrix is possible. None of those
        freedoms moves a leading entry, so the pivot positions are unaffected (see
        Pause and think~\ref{think:pivotpos}). (2)~Consistent, since the last column (column~5) is not a pivot column; there are
        $4-2=2$ free variables ($x_2$ and $x_4$). (3)~False: $x_1+x_2=1,\ x_1+x_2=2$ reduces to
        $\begin{amat}{cc|c}1&1&1\\0&0&1\end{amat}$, which has a free variable $x_2$ and is inconsistent.
\end{description}

\end{document}
